\documentclass[11pt]{article}

\usepackage[margin=1in]{geometry}
\usepackage{booktabs}
\usepackage{amsmath}
\usepackage{hyperref}
\usepackage{xcolor}
\usepackage{array}
\usepackage{longtable}
\usepackage{parskip}
\usepackage{caption}
\usepackage{fancyhdr}
\usepackage{titlesec}
\usepackage[htt]{hyphenat}

\titleformat{\section}{\normalfont\Large\bfseries}{\thesection}{1em}{}
\titleformat{\subsection}{\normalfont\large\bfseries}{\thesubsection}{1em}{}

\pagestyle{fancy}
\fancyhf{}
\fancyhead[L]{Prediction Accuracy Report v5}
\fancyhead[R]{\thepage}
\renewcommand{\headrulewidth}{0.4pt}

\hypersetup{
    colorlinks=true,
    linkcolor=blue,
    urlcolor=blue,
    citecolor=blue,
    pdftitle={Prediction Algorithm Accuracy Report v5 - Minutes-Trend Adjustment},
    pdfauthor={Sports Analytics Project}
}

\newcommand{\code}[1]{\allowbreak\texttt{\small #1}\allowbreak}

\title{\textbf{Prediction Algorithm Accuracy Report — v5} \\ \large A Minutes-Trend Feature That Actually Worked}
\author{Sports Analytics Project}
\date{September 2026}

\begin{document}

\maketitle

\begin{abstract}
\noindent Prior reports (v2, v3) rejected two attempts at a minutes-aware player-point predictor: a
multiply-based approach (predicted minutes $\times$ predicted per-minute rate, replacing the existing
prediction) that compounded two noisy estimates and made RMSE worse in both models it was tried on. This
report tests a third, structurally different idea — a small \emph{adjustment} layered on top of the existing
recency-weighted prediction, driven by whether a player's recent minutes have trended away from their own
longer-run baseline, rather than replacing the prediction with a minutes-derived one. Backtested with the
same discipline as every prior experiment in this series (chronological train/validation split, cross-checked
on two independent splits, a targeted check on players with an actual minutes swing), this one is a real,
consistent win for the fantasy-points predictor (\code{apps/optimizer/predict.py}) and a genuine null result
for the scorer-points predictor (\code{game-detail.service.ts}). Shipped to the former only.
\end{abstract}

\tableofcontents

\section{Why a Third Attempt at Minutes}

Two of the three ideas named as "what's next" in earlier reports (opponent-adjustment, a multiply-based
minutes model) were tried and rejected — see v2 and v3's module-comment write-ups. Both failures shared a
root cause: they replaced an existing single-source-of-noise prediction with a derived quantity built from
two independently noisy estimates, which increases variance even when it nudges the mean estimate in the
right direction. This report deliberately avoids that failure mode: the existing recency-weighted prediction
is kept as the baseline in full, and only a small, capped, damped correction is layered on top of it.

\section{Method}

\subsection{The adjustment}

For each prediction, a \emph{trend ratio} is computed: the mean of a player's most recent 3 games' minutes,
divided by the mean of their most recent 10 games' minutes. A ratio above 1 means recent minutes exceed the
longer-run baseline (a role expansion); below 1 means the opposite (a role reduction, injury return still
ramping up, etc.). The existing recency-weighted prediction is then scaled by
$1 + \text{strength} \times \text{clamp}(\text{ratio} - 1, -1, 1)$ — a damped, capped nudge, not a
replacement. A player with fewer than 10 games of history gets no adjustment (trend is undefined without
enough baseline).

\subsection{Backtest}

The same chronological train/validation methodology as every prior experiment: a grid over the adjustment
strength was searched on a training split, the winner was checked against that split's held-out validation
portion, and the whole process was repeated on a second, independent split (70/30 and 60/40) to confirm the
winner wasn't specific to one split's particular cut point.

\begin{table}[h]
\centering
\small
\begin{tabular}{lrrr}
\toprule
Model & Split & Strength (train-winner) & Validation edge (MAE) \\
\midrule
Fantasy points & 70/30 & 0.15 & +0.0149 \\
Fantasy points & 60/40 & 0.15 & +0.0208 \\
Scorer points & 70/30 & 0.05 & +0.0010 \\
Scorer points & 60/40 & 0.00 & +0.0000 \\
\bottomrule
\end{tabular}
\caption{Grid-search winner and validation-split edge (positive = adjustment helps), both models, both
splits. Fantasy points shows the same winning strength on both splits with a consistent, non-trivial edge.
Scorer points shows a negligible-to-zero edge, and its 60/40 split's own grid search picked strength=0.0 —
no adjustment — as optimal.}
\end{table}

\subsection{Targeted check: does the effect concentrate where it should?}

If the mechanism is real, it should help most specifically on players whose minutes actually changed
recently — not uniformly across the whole population, where most players' minutes are stable and there's
nothing for a trend signal to catch.

\begin{table}[h]
\centering
\small
\begin{tabular}{lrr}
\toprule
Model & Split & Edge on players with $|$trend ratio $-$ 1$|$ $>$ 0.3 \\
\midrule
Fantasy points & 70/30 & +0.0672 \\
Fantasy points & 60/40 & +0.0796 \\
Scorer points & 70/30 & +0.0063 \\
Scorer points & 60/40 & +0.0000 \\
\bottomrule
\end{tabular}
\caption{Edge restricted to the subset of predictions where the player had a genuine recent minutes swing.
Fantasy points' edge is 4-5$\times$ larger on this subset than on the full population — exactly the
concentration pattern a real minutes-driven effect should show. Scorer points shows no such concentration.}
\end{table}

This is a meaningfully stronger form of evidence than the aggregate MAE numbers alone: fantasy points'
edge growing substantially larger specifically on the population the adjustment targets, on both independent
splits, is hard to explain as noise. Scorer points showing no concentration at all — not even a modest one —
is equally informative in the other direction.

\section{Why the Two Models Diverge}

\code{game-detail.service.ts}'s scorer-points predictor already caps its recency window at
\code{MOST\_RECENT\_GAMES\_CONSIDERED} = 10 games, versus \code{predict.py}'s fantasy-points predictor, which
uses an unbounded window (every game in a player's history, geometrically down-weighted). A capped 10-game
window is already substantially more locally responsive to a recent role change than an unbounded one — the
oldest, most stale games simply fall out of the window entirely rather than being merely down-weighted. This
is a plausible, mechanistic explanation for why a separate minutes-trend signal has room to add value in one
predictor and not the other, rather than the two predictors' underlying reality actually differing.

\section{What Shipped}

\code{apps/optimizer/predict.py} now computes a minutes-trend ratio per player (3-game recent window vs.\
10-game baseline window) and applies a damped, capped adjustment (strength 0.15, deviation capped at
$\pm 1.0$) on top of the existing recency-weighted fantasy-point prediction. \code{PlayerPrediction} was
regenerated for all players under the new model. \code{game-detail.service.ts} was left unchanged, with the
null result documented in its module comment alongside the two previously-rejected ideas.

11 new unit tests (\code{test\_predict.py}, in a newly-created \code{apps/optimizer} test suite — this app had
no tests before this change) cover the trend-ratio calculation's edge cases (insufficient history, zero
baseline minutes) and the adjustment function's behavior (no-op at ratio 1.0 or \code{None}, correct
direction and magnitude, clamping at the deviation cap). All pass, alongside the full existing 187-test suite
across the predictor, API, and web apps (198 total after this addition), with no regressions.

\section{Recommendations Going Forward}

\begin{enumerate}
    \item \textbf{This is now the third consecutive experiment that isolated one variable, checked it two
    independent ways, and shipped only what survived both} (multi-season Elo tuning in v3, the season-boundary
    reset in v4, this minutes-trend adjustment in v5) — alongside three rejected ideas in the same period
    (opponent-adjustment, the multiply-based minutes model, and this same trend adjustment on the scorer-points
    predictor specifically). The hit rate is roughly even, which is the expected and healthy outcome of
    actually testing ideas rather than shipping on intuition.
    \item \textbf{Four Factors remains the least-improved model across this entire report series} and has had
    no isolated, backtested experiment run against it directly (only the trade-attribution fix and the passive
    benefit of more data). It is the most under-explored remaining target — a natural next step given the
    other three models have each now had at least one dedicated, validated improvement.
    \item \textbf{The scorer-points predictor's capped 10-game window} being the plausible reason it didn't
    benefit from this adjustment is itself testable, not just a hypothesis — a follow-up could check whether
    widening or removing that cap changes the minutes-trend result, though that would need its own isolated
    backtest against the window-size question first.
\end{enumerate}

\section*{Appendix: Reproducing These Results}

The grid-search and validation-split backtest script behind this report is one-off analysis, not part of the
committed codebase, matching the same appendix note in v3 and v4. \code{test\_predict.py} in
\code{apps/optimizer} exercises the shipped adjustment's exact logic (not an approximation of it) and is the
durable, re-runnable check that the shipped behavior matches what was validated here.

\end{document}
